\documentclass[11pt,a4paper]{article}

%% PACHETE MATEMATICE
\usepackage{amsmath,amssymb,amsfonts}
\usepackage{amsthm}
\usepackage{mathpazo}
\usepackage{eulervm}
\usepackage{enumerate}
\usepackage{color}
\usepackage{graphicx}
\usepackage[all]{xy}
\usepackage{xcolor}
\usepackage{framed}
\usepackage[unicode]{hyperref}
\setcounter{MaxMatrixCols}{20}
\usepackage{multicol}
%\usepackage[romanian]{babel}


%% ASPECT
\linespread{1.05}
\usepackage[T1]{fontenc}
\usepackage[utf8x]{inputenc}
\usepackage[margin=2cm]{geometry}
% \usepackage[color]{showkeys}
\usepackage{anyfontsize}
\usepackage{titlesec}
\titleformat*{\section}{\large\bfseries}

\usepackage{fancyhdr}
\pagestyle{fancy}
\rhead{\small \color{gray} Notițe de Adrian Manea}
\lhead{\small \color{gray} M1-ACS, 2017---2018, M. Olteanu}


\usepackage[autostyle = false, style = german]{csquotes}
\usepackage[romanian]{babel}
\newcommand\qq{\enquote}

%% COMENZILE MELE
\renewcommand*{\proofname}{Dem.:}
\newcommand\bb{\mathbb}
\newcommand\dr{\mathrm}
\newcommand\kal{\mathcal}
\newcommand\fk{\mathfrak}
\newcommand\op{\oplus}
\newcommand\ot{\otimes}
\newcommand\ds{\displaystyle}
\newcommand\id{\indent}
\newcommand\nid{\noindent}
\newcommand\seq{\subseteq}
\newcommand\sneq{\subsetneq}
\newcommand\speq{\supseteq}
\newcommand\spneq{\supsetneq}
\newcommand\rar{\rightarrow}
\newcommand\Rar{\Rightarrow}
\newcommand\xrar{\xrightarrow}
\newcommand\lar{\leftarrow}
\newcommand\Lar{\Leftarrow}
\newcommand\xlar{\xleftarrow}
\newcommand\lrar{\leftrightarrow}
\newcommand\Lrar{\Leftrightarrow}
\newcommand\ideal{\trianglelefteq}
\newcommand\ai{\text{ a.\^{i}. }}
\newcommand\sk{(\textit{Schi\c{t}\u{a}}) }
\newcommand\rhup{\rightharpoonup}
\newcommand\rhdn{\rightharpoondown}
\newcommand\lhup{\leftharpoonup}
\newcommand\lhdn{\leftharpoondown}
\newcommand\vid{\emptyset}
\newcommand\wh{\widehat}
\newcommand\ol{\overline}
\newcommand\NN{\mathbb{N}}
\newcommand\RR{\mathbb{R}}
\newcommand\ZZ{\mathbb{Z}}
\newcommand\QQ{\mathbb{Q}}
\newcommand\CC{\mathbb{C}}

%% STILURILE MELE
\newtheoremstyle{dotless}{}{}{\itshape}{}{\bfseries}{:}{ }{}

\theoremstyle{dotless}
\newtheorem{theorem}{Teorem\u{a}}[section]
\newtheorem{proposition}{Propozi\c{t}ie}[section]
\newtheorem{lemma}{Lem\u{a}}[section]
\newtheorem{corollary}{Corolar}[section]
\newtheorem{exercise}{Exerci\c{t}iu}

\newtheoremstyle{dotlessdr}{}{}{}{}{\bfseries}{:}{ }{}
\theoremstyle{dotlessdr}
\newtheorem{example}{Exemplu}[section]
\newtheorem{definition}{Defini\c{t}ie}[section]
\newtheorem{remark}{Observa\c{t}ie}[section]




\begin{document}

\begin{center}
  {\large\textbf{Seminar 10 \\ Integrale duble și triple}}
\end{center}

\vspace{1cm}

\section{Măsură și integrală Lebesgue}
\label{sec:int-leb}

Ideea integrării în mai multe dimensiuni pornește cu o redefinire a conceptului
de \emph{măsură}. Dacă în plan, integrala Riemann permitea calculul de lungimi,
arii și volume pentru funcții de o variabilă, într-un spațiu multidimensional,
este necesară adaptarea măsurii respective, i.e.\ a funcției cu ajutorul căreia
exprimăm ``cantități măsurabile''. Acestea vor fi generalizări ale conceptelor
de tip \emph{lungime, arie, volum}, în sensul că, în cazul particular al două
sau trei dimensiuni, vor coincide cu intuiția geometrică.

Fără a intra în detalii suplimentare, menționăm că măsura care se utilizează
în acest caz, al problemelor din $\RR^n$, se numește \textbf{măsura Lebesgue}.

În particular, \emph{elementul de lungime} față de măsura Lebesgue, va fi notat
$dl$ sau $dx$, \emph{elementul de arie}, cu $dA = dx dy$, iar elementul de volum,
$dV = dxdydz$ (sau corespunzător unor eventuale schimbări de coordonate).

Față de măsura Lebesgue, putem defini \emph{integrala Lebesgue} din $\RR^k$, pe
o submulțime $A \seq \RR^k$, pe care o vom nota prin:
\[
  \int_A f(x_1, \dots, x_k) dx_1 dx_2 \cdots dx_k.
\]

În cazurile particulare $k = 1, 2, 3$, vom folosi notațiile uzuale:
\[
  \int_A f(x) dx, \quad \iint_A f(x, y) dx dy, \quad \iiint_A f(x, y, z) dx dy dz.
\]

Următorul rezultat ne arată că putem inversa ordinea de integrare:

\begin{theorem}[Fubini]\label{thm:fubini}
  Fie $(x, y) \in \RR^{k + p}$ un punct oarecare, notăm măsura Lebesgue din $\RR^k$
  cu $dx$, măsura Lebesgue din $\RR^p$ cu $dy$, iar măsura Lebesgue din $\RR^{k + p}$
  cu $dx dy$.

  Fie $f : \RR^{k + p} \to \RR$ o funcție integrabilă Lebesgue. Atunci are loc:
  \[
    \int_{\RR^k} \Big( \int_{\RR^p} f(x, y) dy \Big) dx = %
    \iint_{\RR^{k+p}} f(x, y) dx dy = %
    \int_{\RR^p} \Big( \int_{\RR^k} f(x, y) dx \Big) dy.
  \]
\end{theorem}

Următoarele aplicații ale teoremei vor fi utile în exerciții.

(1) Fie $\varphi, \psi : [a, b] \to \RR$ două funcții continue astfel încît
$\varphi \leq \psi$ și fie mulțimea:
\[
  K = \{(x, y) \in \RR^2 \mid x \in [a, b], \varphi(x) \leq y \leq \psi(x) \}.
\]

Fie $f : K \to \RR$ o funcție continuă. Atunci $f$ este integrabilă Lebesgue
pe $K$ și are loc:
\[
  \iint_K f(x, y) dx dy = \int_a^b \Big( \int_{\varphi(x)}^{\psi(x)} f(x, y) dy \Big) dx.
\]

În particular, putem calcula \emph{aria mulțimi} $K$ cu funcția identitate:
\[
  \mu(K) = \iint_K dx dy = \int_a^b (\psi(x) - \varphi(x)) dx.
\]

(2) Fie acum $D \seq \RR^2$ o mulțime compactă și fie $\varphi, \psi : D \to \RR$
două funcții continue, cu $\varphi \leq \psi$. Considerăm mulțimea:
\[
  \Omega = \{ (x, y, z) \in \RR^3 \mid (x, y) \in D, \varphi(x, y) \leq z \leq \psi(x, y) \}.
\]

Fie $f : \Omega \to \RR$ o funcție continuă. Atunci $f$ este integrabilă Lebesgue
pe $\Omega$ și are loc:
\[
  \iiint_\Omega f(x, y, z) dx dy dz = \iint_D \Big( \int_{\varphi(x, y)}^{\psi(x, y)} %
  f(x, y, z) dz \Big) dx dy.
\]

În particular, putem calcula \emph{volumul mulțimii} $\Omega$ cu funcția identitate:
\[
  \mu(\Omega) = \iiint_\Omega dx dy dz = \iint_D (\psi(x, y) - \varphi(x, y)) dx dy.
\]

Atunci cînd vrem să calculăm o integrală cu \emph{metoda schimbării de variabile},
teorema următoare ne arată procedura:

\begin{theorem}[Schimbare de variabile]\label{thm:sch-var}
  Fie $A \seq \RR^n$ o mulțime deschisă și fie $\varphi$ un \emph{difeomorfism} \footnote{%
    aplicație continuă și bijectivă, cu inversa continuă} al lui $A$. Pentru orice
  funcție continuă $f : \varphi(A) \to \RR$, are loc:
  \[
    \int_{\varphi(A)} f(x) dx = \int_A (f \circ \varphi)(y) \cdot |J_\varphi (y)| dy,
  \]
  unde $J_\varphi$ este jacobianul difeomorfismului $\varphi$.
\end{theorem}

%%%%%%%%%%%%%%%%%%%%%%%%%%%%%%%%%%%%%%%%%%%%%%%%%%%%%%%%%%%%%%%%%%%%%%

\section{Exerciții}
\label{sec:exc}

1. Să se calculeze următoarele integrale duble:
\begin{enumerate}[(a)]
\item $\ds\iint_D xy^2 dx dy$, unde $D = [0, 1] \times [2, 3]$;
\item $\ds\iint_D xy dx dy$, unde $D = \{(x, y) \in \RR^2 \mid y \in [0, 1], y^2 \leq x \leq y \}$;
\item $\ds\iint_D y dxdy$, unde $D = \{(x, y) \in \RR^2 \mid (x - 2)^2 + y^2 \leq 1\}$.
\end{enumerate}

\emph{Indicații:} Integralele duble se pot gîndi ca \emph{integrale succesive}, adică:
\[
  \iint_D xy^2 dx dy = \int_0^1 \int_2^3 xy^2 dy dx = \int_0^1 \Big( \int_2^3 xy^2 dy \Big) dx.
\]

\vspace{1cm}

2. Să se calculeze integralele duble:
\begin{enumerate}[(a)]
\item $\ds\iint_D (x + 3y) dx dy$, unde $D$ este mulțimea plană mărginită de
  curbele de ecuații:
  \[
    y = x^2 + 1, \quad y = -x^2, \quad x = -1, \quad x = 3.
  \]
\item $\ds\iint_D e^{|x + y|} dx dy$, unde $D$ este mulțimea plană mărginită de
  curbele de ecuații:
  \[
    x + y = 3, \quad x + y = -3, \quad y = 0, \quad y = 3.
  \]
\item $\ds\iint_D xdxdy$, unde $D$ este mulțimea plană mărginită de:
  \[
    x^2 + y^2 = 9, \quad x \geq 0.
  \]
\end{enumerate}

\emph{Indicații:}
\begin{enumerate}[(a)]
\item $\ds\iint_D (x + 3y) dx dy = \ds\int_{-1}^3 dx \ds\int_{-x^2}^{x^2 + 1} (x + 3y) dy$;
\item Pentru explicitarea modulului, considerăm $D_1 = \{(x, y) \in D \mid x + y \leq 0 \}$
  și $D_2 = D - D_1$. Atunci avem:
  \begin{align*}
    \iint_D e^{|x + y|} dx dy &= \iint_{D_1} e^{-x-y} dx dy + \iint_{D_2} e^{x+y} dx dy \\
                              &= \int_0^3 dy \int_{-3-y}^{-y} e^{-x - y} dx + %
                                \int_0^3 dy \int_{-y}^{3 - y} e^{x + y} dx.
  \end{align*}
\item $\ds\iint_D xdxdy = \int_{-3}^3 dy \int_0^{\sqrt{9 - y^2}} xdx$.
\end{enumerate}

\vspace{1cm}

3. Folosind coordonatele polare, să se calculeze:
\begin{enumerate}[(a)]
\item $\ds\iint_D e^{x^2 + y^2} dx dy$, unde $D = \{(x, y) \in \RR^2 \mid x^2 + y^2 \leq 1 \}$;
\item $\ds\iint_D \big( 1 + \sqrt{x^2 + y^2} \big) dx dy$, pe mulțimea:
  \[
    D = \{(x, y) \in \RR^2 \mid x^2 + y^2 - y \leq 0, x \geq 0\}.
  \]
\item $\ds\iint_D \ln(1 + x^2 + y^2) dx dy$, unde $D$ este mărginit de:
  \[
    x^2 + y^2 = e^2, \quad y = x \sqrt{3}, \quad x = y\sqrt{3}, \quad x \geq 0.
  \]
\end{enumerate}

\emph{Indicație:} Coordonatele polare sînt $x = r \cos \theta, y = r \sin \theta$,
jacobianul este $r$, iar domeniul maxim de definiție este:
\[
  (r, \theta) \in [0, \infty) \times [0, 2\pi).
\]

(b) Obținem restricții suplimentare pentru domeniu, din:
\[
  r \leq \sin \theta, \quad \cos \theta \geq 0 \Rightarrow %
  \theta \in \big[0, \frac{\pi}{2}\big), \quad r \in [0, \sin \theta].
\]

\vspace{1cm}

4. Fie $D \seq \RR^2$ și $f : D \to [0, \infty)$. Să se calculeze volumul
mulțimii:
\[
  \Omega = \{(x, y, z) \in \RR^3 \mid (x, y) \in D, 0 \leq z \leq f(x, y) \},
\]
în următoarele cazuri:
\begin{enumerate}[(a)]
\item $D = \{(x, y) \in \RR^2 \mid x^2 + y^2 \leq 2y \}, f(x, y) = x^2 + y^2$;
\item $D = \{(x, y) \in \RR^2 \mid x^2 + y^2 \leq x, y > 0 \}, f(x, y) = xy$;
\item $D = \{(x, y) \in \RR^2 \mid x^2 + y^2 \leq 2x + 2y - 1 \}, f(x, y) = y$.
\end{enumerate}

\emph{Indicație:} $\mu(\Omega) = \ds\iint_D f(x, y) dx dy$. Folosiți coordonate polare.

\vspace{1cm}

5. Să se calculeze ariile mulțimilor plane $D$ mărginite de curbele de ecuații:
\begin{enumerate}[(a)]
\item $\frac{x^2}{a^2} + \frac{y^2}{b^2} = 1$ (elipsa), cu $a, b > 0$;
\item $(x^2 + y^2)^2 = a^2(x^2 - y^2), x > 0, a > 0$;
\item $(x^2 + y^2)^2 = 2a^2 xy$ (lemniscata lui Bernoulli) \footnote{\color{blue} \href{https://en.wikipedia.org/wiki/Lemniscate_of_Bernoulli}{[wiki]}},
  cu $a > 0$.
\end{enumerate}

\emph{Indicație:} $\mu(D) = \ds\iint_D dx dy$. Folosiți coordonate polare.

\end{document}
